% SPDX-License-Identifier: Apache-2.0
\section{The Lowering, Checked}
\label{sec:x-checked}

A lowering that is only printed is a claim. From this section on, a
lowered unit is also run, and against the same stimulus the Rust
simulation had. The build does it: the example is run with
\code{TXHDL\_VHDL} set, so beside its trace it writes the unit's VHDL,
which \code{\#[lower]} renders from the same statements as the
Verilog, and a ports file; \code{//tools/fst2tb} reads the trace and
the ports and writes a testbench that replays the trace's inputs into
the entity and asserts, tick by tick, that the entity's registers and
outputs are what the trace holds; and \code{vhdl\_test} from
\code{rules\_nvc} simulates it under \code{nvc}, built from source,
hermetically. The test passes only if every comparison held.

The testbench's timing is the model of time, made literal. One tick is
one nanosecond and the clock rises at even ticks. An input the trace
holds at tick $2k$ was set before that edge, so it is applied at
$2k-1$. A register the trace holds at $2k$ took its value at that
edge, so it is checked at $2k+1$, through a VHDL-2008 external name.
A wire the Rust process drove at $2k$ was computed from the registers
as they stood before that edge, and from the inputs of that edge, so
it is checked at $2k-1$ after those inputs are applied and a delta
has passed, where the entity's continuous assignment shows the same
value. That last rule
is a fact about the runtime worth knowing: a wire a process drives
inside its iteration lags the hardware's continuous assignment by the
process's own step.

The counter of the cheat sheet, the blinky, the sequencer, the ALU,
the scratchpad and the decoder pass.
A negative control was run once: the counter's adder changed to add
two makes the test fail at the first tick, so a pass means something.

The stage of Section~\ref{sec:x-stage} passes too, and that closes a
finding. The runtime's channel was a mailbox: an offer sat until
taken, and the consumer took it in the same step it was offered, so
\code{valid} was never high at a tick's end and \code{ready} meant
the slot was empty, while the lowering's channel was a valid/ready
handshake held across edges; a trace of one was not a check of the
other. The runtime's channel is now an elastic buffer of two with the
handshake registered on both sides, as the runtime document states,
and the netlist's port is what that buffer shows: the head and the
room are inputs the testbench applies a step later than a wire, since
they are registered state, and the take and the offer are outputs
checked as any wire is, the offer's data only under its valid. The
change moved the FIFO's, the join's and the arbiter's traces by the
buffer's latency and changed nothing else about them.
